% Source-faithful assembly: Mumford, Abelian Varieties, Appendix II,
% Horizon transcription pages 0282--0286.

\section{Heights associated with line bundles}
\label{mumford-appendix-II-heights-line-bundles}
\dcref{M:appII:II}

For a morphism $\varphi:X\to\mathbb P^n$ over $K$, put
$h_\varphi(x)=h(\varphi(x))$ on $X(\overline K)$.

\begin{proposition}[Height of a line bundle embedding]
\label{mumford-prop-morphism-height-equivalence}
Let $X$ be complete over $K$, and let $\varphi:X\to\mathbb P^k$ and
$\psi:X\to\mathbb P^l$ be $K$-morphisms. If
$\varphi^*\mathcal O(1)\simeq\psi^*\mathcal O(1)$, then
$h_\varphi\sim h_\psi$ on $X(\overline K)$.
\source{mumford-abelian-varieties:page-0282}
\end{proposition}
\begin{proof}
Put $L=\varphi^*\mathcal O(1)=\psi^*\mathcal O(1)$ and let
$\chi:X\to{\rm Proj}(S(\Gamma(X,L)))$ be the associated morphism. Then
$\chi^*\mathcal O(1)=L$, so it suffices by symmetry to prove
$h_\chi\sim h_\varphi$. Assume $\varphi(X)$ is not contained in a proper
linear subspace. Choose $(T_0,\ldots,T_k)$ from the pulled-back coordinate
sections and extend it to $(T_0,\ldots,T_n)$ in $\Gamma(X,L)$. For
$\chi(x)=(x_0,\ldots,x_n)$ and $\varphi(x)=(x_0,\ldots,x_k)$,
$h_\varphi\le h_\chi$.

The image $\chi(X)$ is closed. If $I$ is its homogeneous ideal and
$R=K[T_0,\ldots,T_n]/I$, the first $k+1$ sections have no common zero, so
some $q>0$ and forms $F_{ij}$ of degree $q-1$ satisfy
\[
 T_i^q=\sum_{j=0}^kF_{ij}(T_0,\ldots,T_n)T_j,
 \qquad i=1,\ldots,n-k. \tag{7}
\]
Hence
$x_i^q=\sum_{j=0}^kF_{ij}(x_0,\ldots,x_n)x_j$. Applying the height estimate
gives
\[
 q\log|x_i|_v\le(q-1)\max_{j\le n}\log|x_j|_v+
 \max_{j\le n}\log|x_j|_v+C_v,
\]
where $C_v$ is nonzero at only finitely many places. Thus
$\max_{i\le n}\log|x_i|_v\le\max_{j\le k}\log|x_j|_v+C_v$, and
$h_\chi\le h_\varphi+C$.

\end{proof}


\begin{theorem}[Weil height theorem]
\label{mumford-thm-weil-height-line-bundle}
For every projective variety $X/K$ and $L\in{\rm Pic}X$, there is a height
$h_L$ on $X(\overline K)$, unique up to equivalence, such that
\begin{enumerate}
\item $h_{L_1\otimes L_2}\sim h_{L_1}+h_{L_2}$;
\item for $X=\mathbb P^n$, $h_{\mathcal O(1)}$ is the projective height;
\item for $\varphi:X\to Y$, $h_{\varphi^*L}\sim h_L\circ\varphi$.
\end{enumerate}
\source{mumford-abelian-varieties:page-0283}
\source{mumford-abelian-varieties:page-0284}
\uses{mumford-prop-morphism-height-equivalence,mumford-lem-projective-height}
\end{theorem}
\begin{proof}
Uniqueness follows because ${\rm Pic}X$ is generated by very ample bundles,
whose heights have the required properties by the proposition. For
$L_1=\varphi^*\mathcal O(1)$ and $L_2=\psi^*\mathcal O(1)$, use the Segre map
$\sigma:\mathbb P^k\times\mathbb P^l\to\mathbb P^N$:
\[
 ((x_0,\ldots,x_k),(y_0,\ldots,y_l))\mapsto(\ldots,x_iy_j,\ldots),
\]
with
$\sigma^*\mathcal O(1)=p_1^*\mathcal O(1)\otimes p_2^*\mathcal O(1)$.
The composite $\chi=\sigma\circ(\varphi,\psi)$ represents $L_1\otimes L_2$,
and
\[
 \max_{i,j}\log|x_iy_j|_v=\max_i\log|x_i|_v+\max_j\log|y_j|_v,
\]
so $h_\chi\sim h_\varphi+h_\psi$. For arbitrary $L$, write
$L=L_1\otimes L_2^{-1}$ with $L_i$ very ample and define
$h_L=h_{L_1}-h_{L_2}$. The preceding result makes this independent up to
equivalence and proves (a),(b). For (c), it suffices to use generators
$L=\psi^*\mathcal O(1)$; then both sides are the height of
$\psi\circ\varphi$.

\end{proof}


\begin{lemma}[Approximate quadratic functions]
\label{mumford-lem-approximate-quadratic}
If $h:\Gamma\to\mathbb R$ on an abelian group satisfies
\[
h(x_1+x_2+x_3)-\sum_{i<j}h(x_i+x_j)+\sum_i h(x_i)\sim0,\tag{8}
\]
then there are unique symmetric bilinear $b:\Gamma^2\to\mathbb R$ and
homomorphism $l:\Gamma\to\mathbb R$ such that
\[
h(x)\sim\widehat h(x)=\frac12b(x,x)+l(x).
\]
\source{mumford-abelian-varieties:page-0284}
\source{mumford-abelian-varieties:page-0285}
\end{lemma}
\begin{proof}
Put $\beta(x_1,x_2)=h(x_1+x_2)-h(x_1)-h(x_2)$. It is symmetric and
approximately additive:
\[
\beta(x_1+x_2,x_3)-\beta(x_1,x_3)-\beta(x_2,x_3)\sim0.\tag{9}
\]
Define
\[
b(x_1,x_2)=\lim_{n\to\infty}4^{-n}\beta(2^nx_1,2^nx_2).
\]
Indeed, with $\theta_n\sim0$,
\[
4^{-(n+1)}\beta(2^{n+1}x_1,2^{n+1}x_2)=
4^{-n}\beta(2^nx_1,2^nx_2)+4^{-(n+1)}\theta_n,
\]
so
\[
4^{-N}\beta(2^Nx_1,2^Nx_2)=\beta(x_1,x_2)+
\sum_{n=0}^{N-1}4^{-(n+1)}\theta_n.
\]
Equation (9) gives bilinearity of $b$. Set
$\lambda(x)=h(x)-\frac12b(x,x)$. Then
$\lambda(x_1+x_2)-\lambda(x_1)-\lambda(x_2)\sim0$, and the same averaging gives
\[
l(x)=\lim_n2^{-n}\lambda(2^nx),
\]
a homomorphism. Thus $\widehat h\sim\frac12b+l$, and uniqueness is immediate.

\end{proof}


\begin{lemma}[Tate height]
\label{mumford-lem-tate-height}
\label{mumford-lem-height-parallelogram}
For an abelian variety $X/K$ and $L\in{\rm Pic}X$, the Weil height $h_L$
satisfies (8). Hence $b_L,l_L$ and the Tate height $\widehat h_L$ are uniquely
defined.
\source{mumford-abelian-varieties:page-0285}
\source{mumford-abelian-varieties:page-0286}
\uses{mumford-thm-weil-height-line-bundle,mumford-lem-approximate-quadratic}
\end{lemma}
\begin{proof}
Apply the theorem of the cube to the projections from $X^3$; the resulting
line-bundle identity is
\[
\left(\sum_{i=1}^3p_i\right)^*L\otimes
\bigotimes_{i<j}(p_i+p_j)^*L^{-1}\otimes\bigotimes_{i=1}^3p_i^*L\simeq L.
\]
Taking heights at $(x_1,x_2,x_3)$ and using the Weil height properties yields
(8).

\end{proof}


\begin{remark}[Functoriality and symmetric bundles]
\label{mumford-rem-tate-height-functoriality}
For $\varphi:X\to Y$ and $L\in{\rm Pic}Y$,
$\widehat h_{\varphi^*L}=\widehat h_L\circ\varphi$ and
$b_{\varphi^*L}=b_L\circ\varphi$, $l_{\varphi^*L}=l_L\circ\varphi$.
For $L^{-1}$, $b_{L^{-1}}=-b_L$ and $l_{L^{-1}}=-l_L$. If $L$ is symmetric,
$l_L=0$ and $\widehat h_L(x)=\frac12b_L(x,x)$.
\source{mumford-abelian-varieties:page-0286}
\end{remark}

\begin{proof}[Proof of Proposition 3 and Mordell--Weil]
Choose a symmetric very ample $L$ on $X$ and put
$\langle x,y\rangle=b_L(x,y)$. Then $\langle x,x\rangle\ge0$: otherwise
$\widehat h_L(nx)=\frac12n^2\langle x,x\rangle\to-\infty$, contradicting
$\widehat h_L\sim h_L$ and the existence of a nonnegative representative of
the height class. Proposition 4 implies that the set of $x\in X(K)$ with
$h_L(x)\le c$ is finite, hence the same holds for
$\{x:\langle x,x\rangle\le C\}$. Proposition 2 follows, and the formal
finite-generation criterion proves the Mordell--Weil theorem.
\source{mumford-abelian-varieties:page-0286}
\uses{mumford-lem-tate-height,mumford-prop-bounded-height-projective,mumford-prop-formal-finite-generation}
\end{proof}
